\section{A Field inside the Ring}\label{sec:field}

LaBRADOR distinguishes constraints over $\Rq$ from constraints on constant coefficients, that is, over $\Z_q$. The former are compatible with amortisation, the latter are what the norm proof and most applications produce, and passing from the latter to the former is expensive. This section sets up an intermediate level: a field $\FF$ with $\Z_q\subset\FF\subset\Rq$ that is large enough for random linear combinations with negligible error and small enough for its elements to be cheap to send. The field and its trace are classical objects, and subfields of this kind are used in Hachi~\cite{NOZ26}; we recall them with proofs because the constants matter, and then describe how constraints are aggregated and lifted with them, which is the part that the argument of \cref{sec:tail} relies on.

\subsection{The field}

Let $f\le d/2$ be a power of two, $u=d/(2f)$, and
\[
\eta_k=X^{ku}+X^{-ku}=X^{ku}-X^{d-ku}\in R\quad(1\le k<f),\qquad \eta=\eta_1 .
\]
The element $X^u$ is a primitive $4f$-th root of unity, so $S=\Z[\eta]$ is the ring of integers of the maximal real subfield of the $4f$-th cyclotomic field, with $\Z$-basis $1,\eta_1,\dots,\eta_{f-1}$~\cite[Proposition~2.16]{Was97}.

\begin{lemma}\label{lem:field}
Let $q\equiv5\pmod8$. Then $\FF=\Z_q[\eta]\subset\Rq$ is a field with $q^f$ elements, and $\Rq$ is an $\FF$-vector space with basis $1,X,\dots,X^{d/f-1}$.
\end{lemma}
\begin{proof}
The Galois group of $\Q(\eta)$ over $\Q$ is $(\Z/4f)^\times/\{\pm1\}$. Since $(\Z/4f)^\times=\{\pm1\}\times\langle5\rangle$, this group is cyclic of order $f$ and generated by the class of $5$. Write $q\equiv\pm5^j\pmod{4f}$. If $f\ge2$, reduction modulo $8$ shows that the sign is $+$ and that $j$ is odd, because the powers of $5$ are $1$ or $5$ modulo $8$; hence the class of $q$ generates the group. The Frobenius element at $q$ therefore has order $f$, the prime $q$ is inert in $S$, and $S/qS$ is a field with $q^f$ elements. This holds trivially for $f=1$. The coefficient vectors of $1,\eta_1,\dots,\eta_{f-1}$ are part of a $\Z$-basis of $\Z^d$, so $S\cap qR=qS$ and $S/qS$ embeds into $\Rq$ with image $\Z_q[\eta]$.

For the second claim, $X^{2u}-\eta X^u+1=0$ shows that $1,X,\dots,X^{2u-1}$ generate $R$ as an $S$-module. Since $R$ has rank $d=f\cdot2u$ over $\Z$, they form a basis, and so do their images modulo $q$.
\end{proof}

For $d=64$ we take $f=4$: then $\eta=X^8-X^{56}$, the field $\FF$ has $q^4\approx2^{128}$ elements, an element of $\FF$ is stored in $16$ bytes, and $\Rq$ has dimension $16$ over $\FF$. We write elements of $\FF$ in the basis $1,\eta_1,\dots,\eta_{f-1}$ and let $\tau_0\colon\FF\to\Z_q$ be the coordinate of~$1$.

\subsection{The trace}

Let $H=\{\sigma_j\mid j\equiv\pm1\pmod{4f}\}$, the group of the $d/f$ automorphisms of $R$ that fix $S$, and define
\begin{equation}\label{eq:theta}
\Theta\colon\Rq\to\FF,\qquad\Theta(a)=\frac fd\sum_{\sigma\in H}\sigma(a) .
\end{equation}

\begin{lemma}\label{lem:theta}
\begin{enumerate}[label=(\roman*)]
\item $\Theta$ is $\FF$-linear, and $\Theta(c)=c$ for $c\in\FF$.
\item For $a=\sum_{i<d}a_iX^i$ we have $\Theta(a)=a_0+\frac12\sum_{k=1}^{f-1}(a_{ku}-a_{d-ku})\,\eta_k$.
\item $\ct=\tau_0\circ\Theta$.
\item If $\delta\in\Rq$ is non-zero and $\alpha$ is uniformly random in $\Rq$, then $\Theta(\alpha\delta)$ is uniformly random in $\FF$.
\end{enumerate}
\end{lemma}
\begin{proof}
(i) holds because the elements of $H$ are ring homomorphisms that fix $\FF$. For~(ii), the automorphisms $\sigma_j$ with $j\equiv1\pmod{4f}$ contribute $X^i\sum_{t<u}(X^{4fi})^t$ to the trace of $X^i$, and $X^{4f}$ has multiplicative order $u$; the contribution is $uX^i$ if $u$ divides $i$ and $0$ otherwise. With the automorphisms $j\equiv-1$ we get $\Theta(X^{ku})=\frac12(X^{ku}+X^{-ku})$ and $\Theta(X^i)=0$ if $u\nmid i$. Now $X^{ku}+X^{-ku}$ is $2$ for $k=0$, $\eta_k$ for $1\le k<f$, zero for $k=f$, and $-\eta_{2f-k}$ for $f<k<2f$. (iii) follows from~(ii). For~(iv), the map $\alpha\mapsto\Theta(\alpha\delta)$ is $\FF$-linear. It is not zero: otherwise $\delta_i=\ct(X^{-i}\delta)=\tau_0(\Theta(X^{-i}\delta))=0$ for all $i$. A non-zero linear map onto a one-dimensional space is surjective with fibres of equal size.
\end{proof}

\begin{figure}[t]
\centering
\begin{tikzpicture}[font=\small,>=Stealth,
  box/.style={draw,rounded corners=2pt,minimum height=7.5mm,minimum width=31mm,inner sep=2pt}]
\node[box,fill=cZ!25] (R) at (0,2.7) {$\Rq$\quad$64$ coefficients};
\node[box,fill=cSC!28] (F) at (0,1.35) {$\FF\cong\FF_{q^4}$\quad$4$ coefficients};
\node[box,fill=gray!15] (Z) at (0,0) {$\Z_q$\quad$1$ coefficient};
\draw[->] (R) -- node[right,font=\footnotesize]{$\Theta$} (F);
\draw[->] (F) -- node[right,font=\footnotesize]{$\tau_0$} (Z);
\draw[->,rounded corners=4pt] (R.west) -- ++(-0.7,0) |- node[pos=0.25,left,font=\footnotesize]{$\ct$} (Z.west);
\node[anchor=west,align=left,font=\footnotesize] at (2.3,2.7) {constraints $\ip{\vphi}{\vw}=b$:\\ preserved by challenges $c\in R$ (amortisation)};
\node[anchor=west,align=left,font=\footnotesize] at (2.3,1.35) {constraints $\Theta(\ip{\vphi}{\vw})=b$:\\ preserved by challenges $\rho\in\FF$ (aggregation, sum-check)};
\node[anchor=west,align=left,font=\footnotesize] at (2.3,0) {constraints $\ct(\ip{\vphi}{\vw})=b$:\\ projections, evaluations, coefficients};
\end{tikzpicture}
\caption{The three levels of linear constraints for $d=64$ and $f=4$, and the linear maps between them. An element of $\Rq$ takes $256$ bytes, an element of $\FF$ takes $16$.}
\label{fig:tower}
\end{figure}

\subsection{Constraints over the field}

For a vector $\vw\in\Rq^m$ we consider constraints of three kinds (\cref{fig:tower}): \emph{$R$-constraints} $\ip{\vphi}{\vw}=b$ with $b\in\Rq$, \emph{$\FF$-constraints} $\Theta(\ip{\vphi}{\vw})=b$ with $b\in\FF$, and \emph{$\Z_q$-constraints} $\ct(\ip{\vphi}{\vw})=b$ with $b\in\Z_q$; in all cases $\vphi\in\Rq^m$. A constraint over $\Z_q$ on the coefficient vector of $\vw$ is a $\Z_q$-constraint by~\eqref{eq:ct}. By \cref{lem:theta}, random linear combinations turn any number of $R$-constraints and $\FF$-constraints into a single $\FF$-constraint:

\begin{lemma}[aggregation]\label{lem:agg}
Let $\ip{\vphi_k}{\vw}=b_k$ for $k\le K$ be $R$-constraints and $\Theta(\ip{\vphi'_l}{\vw})=b'_l$ for $l\le L$ be $\FF$-constraints. For $\alpha_1,\dots,\alpha_K\in\Rq$ and $\xi_1,\dots,\xi_L\in\FF$ let
\[
\vPhi=\sum_k\alpha_k\vphi_k+\sum_l\xi_l\vphi'_l,\qquad b=\sum_k\Theta(\alpha_kb_k)+\sum_l\xi_lb'_l .
\]
If $\vw$ satisfies all constraints, then $\Theta(\ip{\vPhi}{\vw})=b$. If $\vw$ violates one of them and the $\alpha_k$ and $\xi_l$ are uniformly random, then $\Theta(\ip{\vPhi}{\vw})=b$ with probability $1/\lvert\FF\rvert$.
\end{lemma}
\begin{proof}
By linearity, $\Theta(\ip{\vPhi}{\vw})-b=\sum_k\Theta\bigl(\alpha_k(\ip{\vphi_k}{\vw}-b_k)\bigr)+\sum_l\xi_l\bigl(\Theta(\ip{\vphi'_l}{\vw})-b'_l\bigr)$. If one of the differences in parentheses is not zero, the corresponding term is uniformly random in $\FF$, by \cref{lem:theta}(iv) for the first sum and because $\FF$ is a field for the second, and it is independent of the other terms.
\end{proof}

\paragraph{Lifting.}
A $\Z_q$-constraint $\ct(\ip{\vpsi}{\vw})=p$ is not preserved by multiplication with elements of $\FF$ or $R$. To make it so, the prover \emph{lifts} it: it sends the element $y=\Theta(\ip{\vpsi}{\vw})$ of $\FF$, of which the verifier already knows the coordinate $\tau_0(y)=p$ by \cref{lem:theta}(iii), and the constraint is replaced by the $\FF$-constraint $\Theta(\ip{\vpsi}{\vw})=y$. This costs $f-1$ elements of $\Z_q$. LaBRADOR lifts to $\Rq$: the prover sends the ring element $\ip{\vpsi}{\vw}$, that is, $d$ elements of $\Z_q$, of which $d-1$ are new to the verifier. Since constraints over $\Z_q$ can be aggregated only with error $1/q$, the aggregation and the lifting are repeated $\kq=4$ times in every level of the recursion, for an error of $q^{-4}\approx2^{-128}$: this costs $1024$ bytes in LaBRADOR and $48$ bytes here.

\begin{remark}\label{rem:field}
Hachi and Akita use subfields of cyclotomic rings and their trace maps to pack field elements into ring elements and to move between rings~\cite{NOZ26,DBK26}. Here the same objects serve the argument itself: the constant-term constraints of the projection are lifted to $\FF$, all constraints of a level are collected in one $\FF$-constraint, and this constraint is reduced by a sum-check whose messages are elements of $\FF$. For this the ring has to contain the field, which excludes moduli that split $X^d+1$ completely. If more than $128$ bits are wanted, $f=8$ gives a field with $q^8$ elements at twice the cost per element (\cref{sec:ablation}).
\end{remark}
